% ---------------------------------------------------------------------
% Parte 3 · WebSockets · Mario · 2:00 (lo edita solo Mario)
% Escribe lo que vas a decir, en tus palabras y en 4 a 6 líneas: es tu
% guion, no una copia de las diapositivas. Cita cada dato (\cite[p.~X]{clave}).
% ---------------------------------------------------------------------
WebSocket se estandarizó en 2011. La idea es abrir una conexión entre el navegador y el servidor y no
cerrarla. Arranca con una petición HTTP que pide cambiar de protocolo, y si el servidor acepta, contesta
con un 101 \cite[secc.~2.4]{Murley2021}.

Con polling el navegador pregunta a cada rato. Según Murley y su equipo, cada ida y vuelta gasta cerca de
590 bytes solo en cabeceras \cite[p.~1197]{Murley2021}. Long polling reduce las preguntas, pero cada
mensaje sigue pagando ese costo.

En WebSocket cada mensaje lleva entre 2 y 14 bytes de cabecera \cite[secc.~2.4]{Murley2021}. En
tradingview una actualización pesa 82 bytes, y los autores calculan que con long polling serían unos 499
\cite[p.~1197]{Murley2021}.
